\section{Introduction}



Vacuum polarization is a phenomenon in quantum electrodynamics (QED) by which the vacuum of a charged field polarizes when subject to external electric fields, effectively behaving like a dielectric with $\epsilon>0$. This mechanism was one of the first QED effects studied \cite{dirac_discussion_1934}, with consequences that followed closely \cite{uehling_polarization_1935,heisenberg_consequences_1936}. This mechanism has been already observed, e.g.~in corrections to atomic levels \cite{lamb_fine_1947} (cf. Ref.~\cite{mohr_qed_1998} for a review), while others e.g.~the Schwinger effect \cite{schwinger_gauge_1951} lay beyond current technological advances (cf. Ref.~\cite{fedotov_advances_2023}). 
Particularly, interest in the Schwinger model has seen an increase in the context of lattice gauge theory, as it is one of the simplest models of a gauge theory that one can simulate, thus becoming a playground for testing novel analytical and numerical methods \cite{banuls_mass_2013,banuls_thermal_2015}. 

Defining vacuum polarization  for general backgrounds\footnote{ And other observables relying on the evaluation of products of fields at coinciding points.} is no easy task. Historically, \cite{dirac_discussion_1934} first defined it by assuming that the short-distance behavior of a class of states when background electric fields are present is the same, modulo smooth coefficients, as in the no background field case. In modern bibliography, these states are called Hadamard, see e.g.~\cite{fewster_necessity_2013}, and are particularly interesting in the context of quantum field theory on curved spacetimes, where the same idea is used to define observables necessary for the study of e.g.~perturbative interacting field theory or semi-classical gravity (cf.~\cite{hollands_quantum_2015,hollands_local_2001} for detailed discussions on the definition on such observables).

Semi-classical gravity, a possible first step towards a quantum theory of gravity, couples classical gravity  to the expectation value of the stress-energy tensor, which as stated above needs to be properly defined. To our knowledge, only perturbative corrections to the metric have been studied  \cite{thompson_quantum-corrected_2024}. In this article we study, as a toy model, the effect of the backreaction of a charged scalar field to a background classical electric field in 1+1 dimensional Minkowski spacetime. This setup was already studied in  \cite{ambjorn_properties_1983}, where it was found that ignoring the backreaction of the scalar field leads to instabilities, and it is claimed that its inclusion avoids them.
However, as pointed out in \cite{schlemmer_current_2015}, the expression for $\rho$ used in \cite{ambjorn_properties_1983} cannot be derived from a gauge invariant prescription. Thus, the question arises: Does using the correct renormalization prescription also avoid these instabilities? We will use the prescription described in \cite{wernersson_vacuum_2021} to implement an iterative method to study the effect of the backreaction of the scalar field. It has been found that this expression also alleviates the above mentioned instabilities.  

The layout of this article is as follows: Section 2 ...

\textbf{Maybe add?}: \cite{costeri_hadamard_2025} discusses hadamard states in the pressence of time-like boundaries.

This article summarizes and extends the results of the M.Sc. thesis of Santiago Sanz-Wuhl, written in 2025 at the Institut für Theoretische Physik, Universität Leipzig, under the supervision of Jochen Zahn.